\paragraph{Traza 2.03 (CVP, dimensión 5).}
El objetivo es $t=(809,994,540,1534,150)$. LLL sube $H$ de 0{,}06 a 0{,}99.
Con la base reducida, Babai da el punto $\widehat v=(806,996,532,1544,153)$, a distancia$^2$ 186
(el redondeo directo da 187), con coeficientes $(-86,-32,-71,45,12)$. La enumeración fija los tres
últimos coeficientes, $(12,45,-71)$, y el centro del siguiente queda en $-32{,}50$, casi a mitad
de camino entre dos enteros. Babai había elegido $-32$, pero al probar $-33$ el último centro pasa
a $-85{,}37$ y con $-85$ sale $v^\star=(807,997,549,1528,155)$, a distancia$^2$ 155. Con ese nuevo
límite se recorren las ramas que quedan (7 nodos en total) sin encontrar nada mejor: 155 es el óptimo.

\paragraph{Traza 3.03 (CVP, dimensión 15).}
LLL baja la norma$^2$ máxima de 710\,154 a 431 y sube $H$ de 0{,}24 a 0{,}96.
Babai obtiene $d^2=8$ y la enumeración lo certifica en 15 nodos, con error
$e=t-v=(0,0,0,1,-1,1,-1,-1,0,0,-1,1,0,1,0)$. En 3.04 (SVP con la misma base), certifica
$\|v\|^2=250$ en 16 nodos.

\subsection{LWE y Ring-LWE}
La Tabla~\ref{tab:lwe} separa los tres métodos: fuerza bruta (L1.01--L2.03), CVP (L2.04--L3.05)
y sistema lineal (Ring-LWE). Palabras Ring-LWE completas:\\
R3.01 \texttt{0100111001010000}; R3.02 \texttt{10100011110000010010000000101100};\\
R3.03 \texttt{0011111100001011111011010001010000010001001001100101000100100111}.

\begin{table}[htb]
\centering
\begin{tabular}{@{}lrrrcclr@{}}
\toprule
Tarea & $n$ & $m$ & $q$ & $c$ & Casos & Resultado & Tiempo (ms)\\
\midrule
L1.01 & 3 & 12 & 281 & 1 & $3^{3}$ & $s=(9,162,198)$ & 0{,}14 \\
L1.02 & 3 & 12 & 379 & 4 & $9^{3}$ & $s=(270,288,152)$ & 6{,}33 \\
L1.03 & 3 & 12 & 281 & 1 & $3^{3}$ & bit 0 & 0{,}13 \\
L1.04 & 3 & 12 & 379 & 4 & $9^{3}$ & bit 0 & 7{,}52 \\
L1.05 & 4 & 16 & 379 & 4 & $9^{4}$ & bit 1 & 82{,}15 \\
L2.01 & 4 & 12 & 463 & 4 & $9^{4}$ & $s=(172,96,358,414)$ & 110{,}12 \\
L2.02 & 5 & 18 & 463 & 4 & $9^{5}$ & $s=(283,111,15,131,443)$ & 1494{,}70 \\
L2.03 & 5 & 15 & 541 & 6 & $13^{5}$ & bit 0 & 5727{,}89 \\
\midrule
L2.03 (CVP) & 5 & 15 & 541 & 6 & -- & bit 0 & 47{,}88 \\
L2.04 & 7 & 25 & 541 & 6 & -- & bit 0 & 309{,}05 \\
L2.05 & 10 & 30 & 541 & 6 & -- & bit 1 & 754{,}27 \\
L3.01 & 7 & 18 & 1223 & 13 & -- & $s=(691,281,601,35,742,795,1015)$ & 120{,}29 \\
L3.02 & 7 & 16 & 1223 & 15 & -- & $s=(1035,694,88,84,498,201,541)$ & 81{,}21 \\
L3.03 & 7 & 16 & 1223 & 15 & -- & bit 1 & 80{,}35 \\
L3.04 & 7 & 18 & 1223 & 15 & -- & bit 0 & 115{,}26 \\
L3.05 & 10 & 30 & 1223 & 15 & -- & bit 1 & 889{,}94 \\
\midrule
R3.01 & 16 & 16 & 1399 & 1 & -- & palabra de 16 bits & 8{,}52 \\
R3.02 & 32 & 32 & 3119 & 1 & -- & palabra de 32 bits & 74{,}60 \\
R3.03 & 64 & 64 & 5903 & 1 & -- & palabra de 64 bits & 805{,}40 \\
\bottomrule
\end{tabular}

\caption{LWE por fuerza bruta (arriba) y como CVP (en medio), y Ring-LWE (abajo). $m$: ecuaciones;
$c$: error máximo encontrado; Casos: $(2c+1)^n$. Tiempo: mediana de 5 repeticiones.}
\label{tab:lwe}
\end{table}

\paragraph{Traza L1.05 (fuerza bruta y descifrado, $n=4$, $m=16$, $q=379$).}
Con $c\le3$ ningún candidato pasa la comprobación. Con $c=4$ aparece $s=(47,82,262,161)$, que deja
en las 16 ecuaciones errores entre $-4$ y 4. Para el criptograma $C=(236,352,5,101)$, $B=115$:
$R=C\cdot s\bmod 379=298$ y $d=115-298=-183$, a 6{,}5 de $-q/2$ y lejos de 0, así que el bit es~1.

\paragraph{Traza L3.01 (CVP, $n=7$, $m=18$, $q=1223$).}
LLL baja la norma$^2$ máxima de 6{,}9 millones a 15\,141. Babai obtiene $d^2=1127$ y la enumeración
lo certifica en 18 nodos. El error centrado $e=(11,-3,12,6,8,13,\dots,-2)$ cumple
$\|e\|^2=1127<8070/4=\lambda_1^2/4$, que asegura la unicidad. El secreto figura en la Tabla~\ref{tab:lwe}.

\begin{samepage}
\paragraph{Traza Ring-LWE R3.01 ($n=16$, $q=1399$).}
Centramos $b=(-69,-64,-33,5,\dots)$. El sistema sin ruido da $(0{,}07;\ 1{,}94;\ 1{,}06;\dots)$;
redondeamos a $s=(0,2,1,2,0,2,1,0,0,0,1,0,1,2,0,0)$ y comprobamos $e\in\{-1,0,1\}$.
Al descifrar, $d=c_2-s\cdot c_1$ empieza por $(-2,-697,8,-1,-692,694,\dots)$; el umbral
$q/4=349{,}75$ da los bits $0,1,0,0,1,1,\dots$
\par\end{samepage}

\subsection{Tiempo frente a talla}
\begin{figure}[H]
\centering
% Tiempo frente a talla: un panel por metodo, en el mismo orden que la memoria (GGH, LWE, Ring-LWE).
% Para que se lea sin interpretar el eje logaritmico, cada punto lleva escrito su tiempo (debajo y a
% la derecha) y entre dos puntos seguidos se indica por cuanto se multiplica (encima de la linea).
% Datos: generar_datos.py y generar_datos_lwe.py (\CurvaGGH, \CurvaLWE, \FactoresLWE, ...).
\pgfplotsset{panel/.style={
  scale only axis,width=3.6cm,height=2.8cm,
  log ticks with fixed point,
  /pgf/number format/use comma,/pgf/number format/1000 sep={},
  grid=major,grid style={gray!20},
  title style={font=\small},
  tick label style={font=\small},label style={font=\small},
  nodes near coords,point meta=explicit symbolic,
  every node near coord/.append style={black,font=\footnotesize,anchor=north west,inner sep=1.5pt},
  enlarge y limits=0.25,clip=false}}
\tikzset{factor/.style={font=\footnotesize,anchor=south east,inner sep=1pt}}
\begin{tikzpicture}
\begin{semilogyaxis}[panel,title={(a) GGH: LLL + enumeración},ylabel={Tiempo (ms)},
  xlabel={Dimensión},xmin=1,xmax=24,xtick={2,7,15},ytick={0.1,1,10}]
\addplot[black,mark=*,mark size=1.6pt,line width=0.7pt] coordinates {\CurvaGGH};
\end{semilogyaxis}
\end{tikzpicture}\hfill
\begin{tikzpicture}
\begin{semilogyaxis}[panel,title={(b) LWE: fuerza bruta ($c=4$)},
  xlabel={Incógnitas},xmin=2.6,xmax=6,xtick={3,4,5},ytick={10,100,1000}]
\addplot[black,mark=*,mark size=1.6pt,line width=0.7pt] coordinates {\CurvaLWE};
\FactoresLWE
\end{semilogyaxis}
\end{tikzpicture}\hfill
\begin{tikzpicture}
\begin{semilogyaxis}[panel,title={(c) Ring-LWE: sistema lineal},
  xlabel={Grado},xmin=8,xmax=100,xtick={16,32,64},ytick={10,100,1000}]
\addplot[black,mark=*,mark size=1.6pt,line width=0.7pt] coordinates {\CurvaRing};
\FactoresRing
\end{semilogyaxis}
\end{tikzpicture}

\caption{Tiempo frente a talla, eje vertical logarítmico. Etiquetas: tiempos y factores de aumento.
GGH: mediana por dimensión.}
\label{fig:tiempos}
\end{figure}
En GGH (Figura~\ref{fig:tiempos}a) pasamos de 0{,}1\,ms en dimensión 2 a unos 20\,ms en 15:
tras LLL apenas hay búsqueda. En LWE (b), con $c=4$ cada incógnita multiplica los candidatos por 9;
los tiempos crecen por \TextoFactoresLWE, aunque también varían las ecuaciones y la parada temprana.
L2.03 tarda más de 5\,s por fuerza bruta y \TiempoCVPLdosTres{} como CVP; las de 10 incógnitas
tardan menos de un segundo. En Ring-LWE (c), doblar el grado multiplica los tiempos por
\TextoFactoresRing, próximos al factor 8 de las operaciones de un sistema cúbico.
